\documentclass[a4paper,twoside,11pt]{mwrep}
\usepackage[utf8]{inputenc}
\usepackage[T1]{fontenc}
\usepackage{polski}
\usepackage{graphicx}
\usepackage{xcolor}
%%%%%%%%%%%%%%%%%%%%%%%%%%%%
\usepackage{hyperref}
\hypersetup{
    colorlinks=true,
    linkcolor=black!30!blue,
    filecolor=black!30!magenta, 
    citecolor=black!30!green,      
    urlcolor=black!30!cyan
}
%%%%%%%%%%%%%%%%%%%%%%%%%%%
% pakiety do kodu źródłowego
\usepackage{listings}
\usepackage{inconsolata}

\lstdefinestyle{JavaStyle}{
    language=Java,
    basicstyle=\ttfamily\footnotesize,
    keywordstyle=\color{blue}\bfseries,
    commentstyle=\color{gray}\itshape,
    stringstyle=\color{orange},
    numbers=left,
    numberstyle=\tiny\color{gray},
    numbersep=8pt,
    frame=single,
    framesep=3pt,
    breaklines=true,
    breakatwhitespace=true,
    tabsize=4,
    showstringspaces=false,
    captionpos=b,
    xleftmargin=15pt,
    framexleftmargin=15pt
}

\lstdefinestyle{YamlStyle}{
    basicstyle=\ttfamily\footnotesize,
    keywordstyle=\color{blue},
    commentstyle=\color{gray}\itshape,
    stringstyle=\color{orange},
    numbers=left,
    numberstyle=\tiny\color{gray},
    numbersep=8pt,
    frame=single,
    framesep=3pt,
    breaklines=true,
    tabsize=2,
    showstringspaces=false,
    captionpos=b,
    xleftmargin=15pt,
    framexleftmargin=15pt
}

%%%%%%%%%%%%%%%%%%%%%%%%%%%
% pakiety matematyczne
\usepackage{amsmath}
\usepackage{amsthm}
\usepackage{amssymb}
%%%%%%%%%%%%%%%%%%%%%%%%%%%%
%%%%%%%%%%%%%%%%%%%%%%%%%%%%
% Ustawienia wymiarów strony
\usepackage{geometry}
 \geometry{
    paper=a4paper,
    inner=30mm,         % Inner margin
    outer=20mm,         % Outer margin
    %bindingoffset=10mm, % Binding offset
    nohead,
    top=20mm,           % Top margin
    bottom=25mm,        % Bottom margin
    }
%%%%%%%%%%%%%%%%%%%%%%%%%%%
% komenda dodająca pustą stronę
\newcommand\blankpage{%
    \null
    \thispagestyle{empty}%
    \addtocounter{page}{-1}%
    \newpage}
%%%%%%%%%%%%%%%%%%%%%%%%%%
% wprowadzenie środowisk dla twierdzeń, lematów, definicji i wnisoków
% Uwaga: można zastosować własne nazwy w miejsce thm, lem, defn i rem
\theoremstyle{plain} \newtheorem{thm}{Twierdzenie}
\theoremstyle{plain} \newtheorem{lem}{Lemat}
\theoremstyle{definition} \newtheorem{defn}{Definicja}
\theoremstyle{remark} \newtheorem*{rem}{Wniosek}
%%%%%%%%%%%%%%%%%%%%%%%%%%%%%%%%%%%%%%%%%%%%%%%%%%%%%%%%%%%%%%%%
% Informacje na stronę tytułową
\newcommand{\autor}{\fontfamily{qhv}\fontshape{n}\selectfont\LARGE\bfseries
Jakub Balicki} % Wpisz swoje imię i nazwisko
%
\newcommand{\album}{\fontfamily{qhv}\fontshape{n}\selectfont
148889} % Wpisz swój numer albumu
%
\newcommand{\pltitle}{\fontfamily{qhv}\fontshape{n}\selectfont\LARGE\bfseries
Projekt i implementacja platformy webowej do gry towarzyskiej Mafia z wykorzystaniem architektury mikro-serwisowej.} % Wpisz tytuł swojej pracy po polsku
%
\newcommand{\engtitle}{\fontfamily{qhv}\fontshape{n}\selectfont\LARGE\bfseries
Design and implementation of a web platform for the social game Mafia using a microservices architecture.} %  Wpisz tytuł swojej pracy po angielsku
%
\newcommand{\level}{\fontfamily{qhv}\fontshape{n}\selectfont\LARGE
inżynierska} % Wybierz odpowiednie określenie
%
\newcommand{\kierunek}{\fontfamily{qhv}\fontshape{n}\selectfont\LARGE
Informatyka}% Wybierz odpowiednie określenie
%
\newcommand{\promotor}{\fontfamily{qhv}\fontshape{n}\selectfont\large\bfseries
dr. Adama Marszałka} % Wpisz odpowiednie dane. Pamiętaj o, że nazwiska w jezyku polskim odmieniają się przez przypadki. Użyj
% dopełniacza.
%
%
\newcommand{\rok}{\fontfamily{qhv}\fontshape{n}\selectfont 2025} % Wpisz rok złożenia pracy

%%%%%%%%%%%%%%%%%%%%%%%%%%%%%%%%%%%%%%%%%%%%%%%%%%%%%%%%%%%%%%%%
%%%%%%%%%%%%%%%%%%%%%%%%%%%%%%%%%%%%%%%%%%%%%%%%%%%%%%%%%%%%%%%%
\begin{document}
% Strona tytułowa
 \begin{titlepage}
 \thispagestyle{empty}
 \vspace*{3ex}
\fontfamily{qhv}\fontshape{n}\selectfont
\hspace{-1.3em}\parbox{0.15\textwidth}{\includegraphics[height=0.15\textwidth]{logoPK.png}}
\parbox[c]{0.7\textwidth}
{\centering
\LARGE
{\bfseries Politechnika Krakowska}\\
{\bfseries im. Tadeusza Kościuszki}\\

\smallskip

\LARGE Wydział Informatyki i Matematyki}\hspace{0.3em}
 \parbox{0.2\textwidth}{\includegraphics[width=0.15\textwidth]{logoIT.png}}%\hspace{0.12em}
\parbox[c]{0.15\textwidth}

\vspace{2ex}

 \vspace{0.13\textheight}

 \center{\autor}
 \smallskip
 \center{\large numer albumu: \album}

\vspace*{3ex}

 \center{\pltitle}
 
 \smallskip

 \center{\engtitle}

 \bigskip

 \center{\LARGE\bfseries praca \level \\[1ex]
\fontfamily{qhv}\fontshape{n}\selectfont\LARGE\bfseries na kierunku
 \kierunek\\[1ex]

\vspace{0.18\textheight}


\hspace{0.4\textwidth} \parbox{0.5\textwidth}{Praca przygotowana pod kierunkiem:
\\
%
\bfseries \promotor}


%\vspace{0.03\textheight}

\vfill

 % wpisać rok
 \center{Kraków, \rok}
 \end{titlepage}
\newpage
\blankpage
%%%%%%%%%%%%%%%%%%%%%%%%%%%%%%%%%%%%%%%%%%%%%%%%%%%%%%%%%%%%%%%%%%%%%%%%%%%%%%
% Spis treści
\pagenumbering{roman}
\tableofcontents
\clearpage
%%%%%%%%%%%%%%%%%%%%%%%%%%%%%%%%%%%%%%%%%%%%%%%%%%%%%%%%%%%%%%%%%%%%%%%%%%%%%% 
%%%%%%%%%%%%%%%%%%%%%%%%%%%%%%%%%%%%%%%%%%%%%%%%%%%%%%%%%%%%%%%%%%%%%%%%%%%%%%
% Początek pracy
\pagenumbering{arabic}
\chapter*{Wstęp}
\addcontentsline{toc}{chapter}{Wstęp}
\setcounter{page}{1}

Gra towarzyska \textit{Mafia} powstała w 1986 roku na Wydziale Psychologii Uniwersytetu Moskiewskiego. Jej twórcą był Dmitrij Dawidow, który opracował mechanikę opartą na konflikcie pomiędzy poinformowaną mniejszością (Mafią) a niepoinformowaną większością (Obywatelami). Rozgrywka szybko zyskała popularność najpierw w środowiskach akademickich ZSRR, a następnie rozprzestrzeniła się na cały świat. Według moskiewskiej legendy mechanika gry była wykorzystywana przy szkoleniu dyplomatów i agentów wywiadu, choć te doniesienia pozostają w sferze anegdot.

Przebieg klasycznej rozgrywki opiera się na cyklicznym powtarzaniu dwóch faz: Nocy, podczas której Mafia w tajemnicy wybiera ofiarę, oraz Dnia, gdy wszyscy gracze dyskutują i głosują nad eliminacją podejrzanego. Zwycięstwo przypada tej frakcji, której członkowie przetrwają do końca gry. Szczegółowy opis tradycyjnych reguł znajduje się w Rozdziale~\ref{chap:Definicja_celu_i_zakresu_pracy}.

Pomimo niemal czterdziestoletniej historii i ogromnej popularności, tradycyjna \textit{Mafia} wymaga obecności prowadzącego, nazywanego Mistrzem Gry. Osoba ta odpowiada za zarządzanie przebiegiem rozgrywki: kontroluje czas poszczególnych faz, zbiera tajne głosy w nocy, ogłasza wyniki eliminacji i pilnuje przestrzegania zasad. Problem polega na tym, że prowadzący nie może jednocześnie uczestniczyć w grze jako pełnoprawny gracz, co w małych grupach stanowi istotne ograniczenie.

Celem niniejszej pracy inżynierskiej jest zaprojektowanie i implementacja platformy webowej \textit{eMafia}, która przejmuje wszystkie obowiązki Mistrza Gry. Aplikacja umożliwia rozgrywkę w trybie lokalnym (twarzą w twarz), gdzie każdy uczestnik korzysta z własnego urządzenia mobilnego lub komputera jako interfejsu do otrzymywania informacji o przydzielonej roli, oddawania głosów i śledzenia przebiegu gry. Dzięki temu wszyscy obecni mogą aktywnie uczestniczyć w rozgrywce, a bezpośrednia interakcja między graczami pozostaje zachowana.

Analiza dostępnych rozwiązań, przeprowadzona w Rozdziale~\ref{chap:Analiza_istniejących_rozwiązań}, wykazała brak platform przystosowanych do hybrydowego trybu gry, łączącego fizyczną obecność graczy z pełną automatyzacją logiki rozgrywki. Istniejące aplikacje koncentrują się albo na rozgrywce wyłącznie online, albo stanowią jedynie cyfrowe odwzorowanie fizycznych kart bez wsparcia dla zarządzania przebiegiem gry.

Z punktu widzenia technicznego realizacja założonego celu wymagała zastosowania architektury mikroserwisowej. System został podzielony na trzy niezależne usługi: \texttt{Auth Service} odpowiedzialny za uwierzytelnianie i zarządzanie tokenami JWT, \texttt{Lobby Service} zarządzający pokojami gier i relacjami graczy, oraz \texttt{Game Engine Service} implementujący całą logikę rozgrywki przy użyciu wzorca maszyny stanów. Backend aplikacji zbudowano w oparciu o framework Spring Boot 3.3, wykorzystując język Java i ORM Hibernate do komunikacji z bazą danych PostgreSQL. Warstwa frontendowa została zaimplementowana przy użyciu biblioteki React. Komunikacja w czasie rzeczywistym odbywa się za pośrednictwem protokołu WebSocket z wykorzystaniem biblioteki STOMP, natomiast komunikacja międzyserwisowa wykorzystuje broker wiadomości RabbitMQ. Infrastruktura aplikacji została w pełni skonteneryzowana za pomocą Docker i Docker Compose.

W kolejnych rozdziałach pracy omówiono: definicję celu i zakresu projektu, analizę istniejących rozwiązań, wymagania i architekturę systemu, implementację poszczególnych mikroserwisów, szczegóły realizacji logiki gry oraz wdrożenie i testowanie aplikacji. Rezultatem pracy jest działająca platforma webowa stanowiąca autorski wkład inżynierski w dziedzinę automatyzacji gier towarzyskich.

\chapter{Definicja celu i zakresu pracy}
\label{chap:Definicja_celu_i_zakresu_pracy}

Niniejszy rozdział zawiera analizę mechaniki tradycyjnej gry towarzyskiej \textit{Mafia} oraz identyfikację problemów związanych z jej prowadzeniem w formie klasycznej. Na podstawie przeprowadzonej analizy sformułowano cel pracy inżynierskiej i określono zakres realizowanych funkcjonalności.

\section{Opis tradycyjnej gry Mafia}
\label{sec:Opis_tradycyjnej_gry_Mafia}

\subsection{Zasady rozgrywki}

Gra \textit{Mafia} należy do kategorii gier towarzyskich opartych na mechanice ukrytych ról i dedukcji społecznej. Rozgrywka wymaga udziału minimum pięciu osób, przy czym optymalna liczba graczy mieści się w przedziale od siedmiu do piętnastu uczestników. Przed rozpoczęciem gry każdy gracz otrzymuje w tajemnicy przydzieloną rolę: członka Mafii lub Obywatela. Proporcje między frakcjami zależą od liczby uczestników -- standardowo Mafia stanowi około jednej czwartej składu.

Rozgrywka przebiega w cyklach składających się z dwóch faz:

\begin{enumerate}
    \item \textbf{Faza Nocy} -- gracze zamykają oczy lub opuszczają wzrok. Mistrz Gry budzi członków Mafii, którzy w milczeniu wskazują ofiarę do eliminacji. Po dokonaniu wyboru Mafia ponownie zasypia.
    
    \item \textbf{Faza Dnia} -- wszyscy gracze otwierają oczy. Mistrz Gry ogłasza, kto został wyeliminowany w nocy. Następuje dyskusja, podczas której gracze próbują zidentyfikować członków Mafii. Na zakończenie fazy odbywa się jawne głosowanie nad eliminacją podejrzanego gracza.
\end{enumerate}

Cykl powtarza się do momentu spełnienia jednego z warunków zwycięstwa:
\begin{itemize}
    \item Obywatele wygrywają, gdy wszyscy członkowie Mafii zostaną wyeliminowani.
    \item Mafia wygrywa, gdy liczba żywych członków Mafii zrówna się z liczbą żywych Obywateli lub ją przewyższy.
\end{itemize}

\subsection{Rola Mistrza Gry}

Prowadzenie rozgrywki wymaga obecności Mistrza Gry, który pełni funkcję arbitra i narratora. Do jego obowiązków należy:

\begin{itemize}
    \item losowe przydzielanie ról przed rozpoczęciem gry,
    \item kontrolowanie przebiegu faz nocnych i dziennych,
    \item zbieranie tajnych głosów Mafii podczas nocy,
    \item ogłaszanie wyników eliminacji,
    \item pilnowanie przestrzegania zasad i limitów czasowych,
    \item rozstrzyganie sytuacji spornych.
\end{itemize}

Obecność prowadzącego generuje istotne ograniczenie: osoba pełniąca tę funkcję nie może uczestniczyć w grze jako pełnoprawny gracz. W przypadku małych grup oznacza to utratę znaczącego odsetka potencjalnych uczestników.

\subsection{Warianty i modyfikacje zasad}

W celu zwiększenia różnorodności rozgrywki stosuje się rozszerzenia podstawowych zasad. Najpopularniejsze modyfikacje obejmują wprowadzenie dodatkowych ról o specjalnych zdolnościach:

\begin{itemize}
    \item \textbf{Detektyw} -- podczas nocy może sprawdzić przynależność frakcyjną wybranego gracza,
    \item \textbf{Lekarz} -- może ochronić wybraną osobę przed eliminacją przez Mafię,
    \item \textbf{Błazen} -- wygrywa, jeśli zostanie wyeliminowany podczas dziennego głosowania.
\end{itemize}

Wprowadzenie ról specjalnych znacząco komplikuje prowadzenie gry. Mistrz Gry musi pamiętać o kolejności budzenia poszczególnych ról w nocy, śledzić efekty ich zdolności i prawidłowo rozstrzygać konflikty między akcjami.

\section{Popularyzacja gry w mediach}

Ostatnie lata przyniosły znaczący wzrost popularności formatów wideo opartych na grze \textit{Mafia}. Na polskim YouTube szczególną popularność zyskały produkcje takie jak \textit{Mafia Friza} czy \textit{Sabotażysta} organizowany przez Reziego. Formaty te przyciągają milionowe widownie i demonstrują, że mechanika ukrytych ról w połączeniu z elementami rozrywkowymi stanowi atrakcyjny materiał medialny.

W formatach medialnych Mistrz Gry pełni rolę aktywnego reżysera, dynamicznie balansującego przebieg wydarzeń poprzez wprowadzanie kart specjalnych lub niespodziewanych zwrotów akcji. Sukces komercyjny tych produkcji wskazuje na istnienie rynkowej niszy dla aplikacji automatyzujących prowadzenie gry \textit{Mafia}.
\section{Potrzeba automatyzacji rozgrywki}

Nadrzędną potrzebą, na którą odpowiada niniejsza praca, jest stworzenie platformy webowej automatyzującej obsługę gry w trybie lokalnym. Aplikacja przejmuje wszystkie funkcje administracyjne Mistrza Gry, stając się autonomicznym silnikiem logiki rozgrywki.

Z perspektywy technicznej rozwiązanie problemu Mistrza Gry wymaga spełnienia następujących założeń:

\begin{enumerate}
    \item \textbf{Synchronizacja w czasie rzeczywistym} -- system musi natychmiast dystrybuować aktualizacje stanu gry do wszystkich graczy za pomocą technologii komunikacji dwukierunkowej (WebSocket). Każda zmiana fazy, wynik głosowania czy eliminacja gracza musi być widoczna dla wszystkich uczestników w ułamku sekundy.
    
    \item \textbf{Skalowalność} -- architektura musi umożliwiać jednoczesne prowadzenie wielu niezależnych rozgrywek bez wzajemnego wpływu na wydajność i stabilność systemu.
    
    \item \textbf{Niezawodność} -- system musi poprawnie obsługiwać sytuacje awaryjne, takie jak utrata połączenia przez gracza czy przekroczenie limitu czasu na głosowanie.
\end{enumerate}
\section{Definicja celu pracy}

Celem niniejszej pracy inżynierskiej jest zaprojektowanie i implementacja rozproszonego systemu informatycznego realizującego logikę gry towarzyskiej \textit{Mafia} w środowisku przeglądarkowym. System ma umożliwić przeprowadzenie pełnej rozgrywki bez udziału Mistrza Gry, automatyzując wszystkie mechanizmy związane z zarządzaniem fazami, głosowaniami i stanem graczy.

Realizacja celu obejmuje:
\begin{itemize}
    \item Opracowanie architektury mikroserwisowej z wydzielonymi domenami odpowiedzialności (uwierzytelnianie, zarządzanie pokojami, logika gry, komunikacja).
    \item Implementację mechanizmu maszyny stanowej (\textit{State Machine}) kontrolującej przepływ faz rozgrywki.
    \item Wdrożenie systemu komunikacji w czasie rzeczywistym opartego na protokole WebSocket z warstwą STOMP.
    \item Konteneryzację wszystkich komponentów systemu przy użyciu Docker i orkiestrację za pomocą Docker Compose.
\end{itemize}
\section{Definicja zakresu pracy}

Zakres pracy obejmuje pełny cykl wytwórczy oprogramowania -- od analizy wymagań, przez projekt architektury, aż po implementację działającego prototypu. Poniżej przedstawiono szczegółowy podział zakresu na poszczególne obszary:

\subsection{Zakres funkcjonalny}
\begin{itemize}
    \item System uwierzytelniania i autoryzacji oparty na tokenach JWT z mechanizmem odświeżania sesji.
    \item Moduł zarządzania pokojami gry -- tworzenie, wyszukiwanie, dołączanie przez kod dostępu.
    \item Silnik logiki gry implementujący fazy: dzień (dyskusja i głosowanie) oraz noc (akcje ról specjalnych).
    \item System głosowania z obsługą remisów i automatycznym zliczaniem głosów.
    \item Mechanizm przydzielania ról (Mafia, Obywatel, role specjalne) z zachowaniem losowości i balansu.
    \item Panel administracyjny umożliwiający zarządzanie użytkownikami i monitorowanie aktywnych rozgrywek.
\end{itemize}

\subsection{Zakres technologiczny}
\begin{itemize}
    \item Backend: Java 21 z frameworkiem Spring Boot 3.3, Spring Security, Spring Data JPA.
    \item Frontend: React 18 z wykorzystaniem CSS Modules i biblioteki Axios.
    \item Baza danych: PostgreSQL 15 w kontenerze Docker.
    \item Komunikacja: REST API dla operacji CRUD, WebSocket (STOMP over SockJS) dla komunikacji w czasie rzeczywistym.
    \item Broker wiadomości: RabbitMQ 3.13 jako warstwa pośrednicząca w architekturze event-driven.
    \item Konteneryzacja: Docker i Docker Compose do orkiestracji środowiska deweloperskiego i produkcyjnego.
\end{itemize}

\subsection{Ograniczenia zakresu}
Praca nie obejmuje:
\begin{itemize}
    \item Wdrożenia produkcyjnego w środowisku chmurowym.
    \item Implementacji aplikacji mobilnej (natywnej lub hybrydowej).
    \item Integracji z zewnętrznymi systemami płatności lub reklam.
    \item Zaawansowanych mechanizmów antycheatowych i systemów reputacji.
\end{itemize}
\chapter{Analiza istniejących rozwiązań}
\label{chap:analiza_istniejacych_rozwiazan}

Przed przystąpieniem do projektowania systemu przeprowadzono analizę porównawczą dostępnych rozwiązań umożliwiających prowadzenie gry \textit{Mafia}. Celem analizy było zidentyfikowanie luki technologicznej, którą niniejsza praca ma wypełnić.

\section{Tradycyjna rozgrywka z Mistrzem Gry}

Klasyczne podejście do gry \textit{Mafia} opiera się na obecności Mistrza Gry -- osoby zarządzającej przebiegiem rozgrywki. Format ten, choć najpopularniejszy w środowiskach lokalnych, jest obarczony istotnymi ograniczeniami technicznymi i organizacyjnymi.

\subsection{Obciążenie prowadzącego}

Mistrz Gry musi jednocześnie realizować wiele zadań administracyjnych:
\begin{itemize}
    \item Zarządzanie czasem poszczególnych faz (dzień, noc, głosowanie).
    \item Dystrybucja i weryfikacja przydzielonych ról.
    \item Zliczanie głosów i rozstrzyganie sytuacji spornych (remisy).
    \item Kontrolowanie kolejności wypowiedzi i eliminacja graczy.
    \item Utrzymywanie poufności informacji o rolach.
\end{itemize}

Złożoność tych obowiązków rośnie proporcjonalnie do liczby graczy i wprowadzanych wariantów reguł. W przypadku rozgrywek z większą liczbą ról specjalnych (Detektyw, Lekarz, Mafia z różnymi zdolnościami) prowadzący staje się wąskim gardłem całego procesu.

\subsection{Brak skalowalności}

Format tradycyjny nie pozwala na równoległe prowadzenie wielu rozgrywek bez proporcjonalnego zwiększenia liczby Mistrzów Gry. Każda sesja wymaga dedykowanej osoby zarządzającej, co stanowi barierę organizacyjną przy większych wydarzeniach.

\subsection{Podatność na błędy}

Ręczne zarządzanie stanem gry prowadzi do nieuniknionych pomyłek: błędnego zliczenia głosów, pominięcia akcji roli specjalnej czy przedwczesnego ujawnienia informacji. Błędy te mogą całkowicie zrujnować przebieg rozgrywki i doświadczenie uczestników.

\section{Platformy cyfrowe do gry online}

Drugą kategorię stanowią w pełni zautomatyzowane platformy cyfrowe, których reprezentatywnym przykładem jest gra \textit{Among Us} wydana przez studio InnerSloth.

\subsection{Zalety automatyzacji}

Platformy cyfrowe oferują szereg udogodnień technicznych:
\begin{itemize}
    \item Pełna automatyzacja logiki gry -- system samodzielnie zarządza fazami, głosowaniami i warunkami zwycięstwa.
    \item Wysoka skalowalność -- architektura serwerowa pozwala na jednoczesne prowadzenie tysięcy niezależnych rozgrywek.
    \item Spójność reguł -- eliminacja błędów ludzkich w interpretacji zasad.
    \item Wbudowane mechanizmy komunikacji (głosowe lub tekstowe).
\end{itemize}

\subsection{Ograniczenia platform online}

Pomimo zalet technicznych, platformy takie jak \textit{Among Us} posiadają fundamentalne ograniczenie -- są zoptymalizowane wyłącznie pod kątem interakcji wirtualnej. Interfejs użytkownika, mechanizmy komunikacji i model rozgrywki zakładają fizyczną separację graczy.

W kontekście spotkań lokalnych (IRL -- \textit{In Real Life}) platformy te tracą swoją przewagę:
\begin{itemize}
    \item Eliminują bezpośrednią interakcję twarzą w twarz, która stanowi istotę gry \textit{Mafia}.
    \item Wymuszają komunikację przez urządzenie zamiast naturalnej dyskusji.
    \item Nie wykorzystują możliwości obserwacji mowy ciała i reakcji emocjonalnych uczestników.
\end{itemize}

\section{Identyfikacja luki technologicznej}

Analiza obu podejść ujawnia niezagospodarowaną niszę rynkową i technologiczną:

\begin{table}[h]
\centering
\begin{tabular}{|l|c|c|}
\hline
\textbf{Cecha} & \textbf{Tradycyjna} & \textbf{Platformy online} \\
\hline
Automatyzacja logiki & Brak & Pełna \\
Skalowalność & Niska & Wysoka \\
Interakcja bezpośrednia & Zachowana & Utracona \\
Wymóg prowadzącego & Tak & Nie \\
Optymalizacja dla IRL & Tak & Nie \\
\hline
\end{tabular}
\caption{Porównanie cech istniejących rozwiązań}
\label{tab:porownanie_rozwiazan}
\end{table}

Żadne z powszechnie dostępnych rozwiązań nie łączy automatyzacji logiki gry z zachowaniem bezpośredniej interakcji między uczestnikami. Ta luka stanowi podstawę dla projektu \textit{eMafia}.

\section{Platforma eMafia jako odpowiedź na zidentyfikowane problemy}

Projekt \textit{eMafia} został zaprojektowany jako hybrydowa platforma webowa, która:

\begin{enumerate}
    \item \textbf{Automatyzuje funkcje Mistrza Gry} -- system przejmuje zarządzanie fazami, głosowaniami, eliminacjami i weryfikacją warunków zwycięstwa.
    
    \item \textbf{Zachowuje interakcję bezpośrednią} -- uczestnicy spotykają się fizycznie, a urządzenia (telefony, tablety) służą jedynie jako interfejs do otrzymywania informacji o roli i oddawania głosów.
    
    \item \textbf{Zapewnia skalowalność} -- architektura systemu pozwala na jednoczesne prowadzenie wielu niezależnych rozgrywek.
    
    \item \textbf{Eliminuje błędy ludzkie} -- automatyczne zliczanie głosów, zarządzanie czasem i weryfikacja stanu graczy.
\end{enumerate}

Realizacja tych założeń wymagała zastosowania architektury rozproszonej, komunikacji w czasie rzeczywistym (WebSocket) oraz mechanizmów synchronizacji stanu między wieloma klientami.
\chapter{Analiza wymagań i projektowanie systemu}
\label{chap:analiza_wymagan}

Niniejszy rozdział przedstawia proces analizy wymagań oraz projektowania architektury systemu \textit{eMafia}. Przyjęte rozwiązania architektoniczne wynikają bezpośrednio z natury problemu -- konieczności synchronizacji stanu gry w czasie rzeczywistym między wieloma klientami.

\section{Wymagania funkcjonalne}

Wymagania funkcjonalne definiują zakres możliwości systemu z perspektywy użytkownika końcowego.

\subsection{Zarządzanie użytkownikami}

System musi zapewniać pełen cykl zarządzania kontami użytkowników:
\begin{itemize}
    \item Rejestracja nowego konta z walidacją unikalności nazwy użytkownika i adresu email.
    \item Logowanie z wykorzystaniem poświadczeń (email i hasło).
    \item Wylogowanie z unieważnieniem aktywnej sesji.
    \item Edycja profilu użytkownika (zmiana nazwy wyświetlanej, hasła).
    \item Panel administracyjny do zarządzania kontami (dla użytkowników z uprawnieniami administratora).
\end{itemize}

\subsection{Zarządzanie pokojami gry}

Moduł lobby musi umożliwiać:
\begin{itemize}
    \item Tworzenie nowego pokoju gry z konfigurowalnymi parametrami (maksymalna liczba graczy, czas faz).
    \item Generowanie unikalnego kodu dostępu do pokoju.
    \item Wyszukiwanie publicznych pokojów według nazwy lub kodu.
    \item Dołączanie do istniejącego pokoju przez kod lub wybór z listy.
    \item Wyświetlanie listy graczy oczekujących w pokoju.
    \item Rozpoczęcie gry przez hosta po zebraniu minimalnej liczby uczestników.
\end{itemize}

\subsection{Logika rozgrywki}

Rdzeń systemu -- silnik gry -- musi realizować następujące funkcje:
\begin{itemize}
    \item Losowe przydzielanie ról (Mafia, Obywatel) z zachowaniem określonych proporcji.
    \item Automatyczne zarządzanie cyklem faz: noc (akcje Mafii) $\rightarrow$ dzień (dyskusja) $\rightarrow$ głosowanie $\rightarrow$ wynik.
    \item Obsługa głosowania dziennego z automatycznym zliczaniem głosów.
    \item Obsługa głosowania nocnego (wybór ofiary przez Mafię).
    \item Eliminacja graczy na podstawie wyników głosowań.
    \item Weryfikacja warunków zwycięstwa po każdej eliminacji.
    \item Dystrybucja informacji o stanie gry do wszystkich podłączonych klientów.
\end{itemize}

\subsection{Warunki zwycięstwa}

System musi automatycznie wykrywać zakończenie gry w dwóch scenariuszach:
\begin{enumerate}
    \item \textbf{Zwycięstwo Mafii} -- liczba żywych członków Mafii jest równa lub większa od liczby żywych Obywateli.
    \item \textbf{Zwycięstwo Obywateli} -- wszyscy członkowie Mafii zostali wyeliminowani.
\end{enumerate}

\section{Wymagania niefunkcjonalne}

Wymagania niefunkcjonalne określają atrybuty jakościowe systemu, które determinują wybór technologii i architektury.

\subsection{Komunikacja w czasie rzeczywistym}

Ze względu na naturę gry \textit{Mafia}, system musi zapewniać natychmiastową propagację zmian stanu do wszystkich klientów. Opóźnienie przekraczające 500ms byłoby zauważalne i negatywnie wpływałoby na doświadczenie użytkownika. Wymóg ten eliminuje tradycyjny model request-response na rzecz komunikacji dwukierunkowej (WebSocket).

\subsection{Spójność stanu}

W systemie rozproszonym, gdzie wielu klientów jednocześnie odczytuje i modyfikuje stan gry, kluczowe jest zapewnienie spójności danych. System musi gwarantować, że:
\begin{itemize}
    \item Wszyscy gracze widzą ten sam stan gry w danym momencie.
    \item Operacje modyfikujące stan (głosowanie, eliminacja) są atomowe.
    \item Konflikty wynikające z równoczesnych akcji są rozwiązywane deterministycznie.
\end{itemize}

\subsection{Bezpieczeństwo}

System musi chronić integralność rozgrywki poprzez:
\begin{itemize}
    \item Uwierzytelnianie użytkowników przed dostępem do zasobów.
    \item Autoryzację operacji -- weryfikację uprawnień do wykonania danej akcji.
    \item Ochronę poufności ról -- uniemożliwienie graczom poznania ról innych uczestników.
    \item Walidację wszystkich danych wejściowych po stronie serwera.
\end{itemize}

\subsection{Skalowalność}

Architektura musi umożliwiać obsługę wielu równoczesnych rozgrywek bez degradacji wydajności. Izolacja stanu poszczególnych pokojów gry jest warunkiem koniecznym dla zapewnienia niezależności sesji.

\section{Projekt architektury systemu}

Na podstawie zebranych wymagań zaprojektowano architekturę opartą na wzorcu klient-serwer z warstwą komunikacji w czasie rzeczywistym.

\subsection{Architektura logiczna}

System został podzielony na trzy główne warstwy odpowiedzialności:

\begin{enumerate}
    \item \textbf{Warstwa prezentacji (Frontend)} -- aplikacja webowa React odpowiedzialna za interfejs użytkownika i obsługę interakcji.
    
    \item \textbf{Warstwa aplikacji (Backend)} -- serwer Spring Boot realizujący logikę biznesową, zarządzanie stanem i komunikację.
    
    \item \textbf{Warstwa persystencji} -- baza danych PostgreSQL przechowująca trwałe dane (użytkownicy, pokoje, sesje).
\end{enumerate}

\subsection{Wzorce komunikacji}

W systemie zastosowano dwa komplementarne wzorce komunikacji:

\begin{itemize}
    \item \textbf{REST API} -- dla operacji niewymagających natychmiastowej odpowiedzi: rejestracja, logowanie, tworzenie pokoju, pobieranie list.
    
    \item \textbf{WebSocket z protokołem STOMP} -- dla komunikacji w czasie rzeczywistym: aktualizacje stanu gry, powiadomienia o akcjach graczy, synchronizacja faz.
\end{itemize}

Protokół STOMP (\textit{Simple Text Oriented Messaging Protocol}) został wybrany ze względu na wbudowane wsparcie w ekosystemie Spring oraz możliwość definiowania kanałów subskrypcji (topics).

\subsection{Model domenowy}

Kluczowe encje systemu i ich relacje:

\begin{itemize}
    \item \texttt{User} -- reprezentuje zarejestrowanego użytkownika (id, username, email, password\_hash, role).
    
    \item \texttt{GameRoom} -- reprezentuje pokój gry (id, name, room\_code, status, host\_id, max\_players, settings).
    
    \item \texttt{PlayerInRoom} -- relacja między użytkownikiem a pokojem z dodatkowymi atrybutami sesji (user\_id, room\_id, game\_role, is\_alive, joined\_at).
    
    \item \texttt{RefreshToken} -- token odświeżający dla mechanizmu JWT (id, token, user\_id, expiry\_date, revoked).
\end{itemize}

\subsection{Maszyna stanowa rozgrywki}

Cykl życia rozgrywki został zamodelowany jako maszyna stanowa z następującymi stanami:

\begin{enumerate}
    \item \texttt{LOBBY} -- pokój oczekuje na graczy, host może rozpocząć grę.
    \item \texttt{NIGHT\_VOTING} -- faza nocna, Mafia wybiera ofiarę.
    \item \texttt{NIGHT\_RESULT} -- prezentacja wyniku nocy (eliminacja lub brak).
    \item \texttt{DAY\_VOTING} -- faza dzienna, wszyscy żywi gracze głosują.
    \item \texttt{DAY\_RESULT} -- prezentacja wyniku głosowania dziennego.
    \item \texttt{GAME\_OVER} -- gra zakończona, wyświetlenie zwycięzców.
\end{enumerate}

Przejścia między stanami są wyzwalane przez zdarzenia systemowe (upływ czasu, zebranie wszystkich głosów) lub akcje użytkownika (rozpoczęcie gry przez hosta).

\section{Projekt schematu bazy danych}

Schemat bazy danych PostgreSQL został zaprojektowany z uwzględnieniem normalizacji i integralności referencyjnej.

\subsection{Tabela users}

Przechowuje dane uwierzytelniające i profilowe użytkowników:
\begin{itemize}
    \item \texttt{id} (UUID, PK) -- unikalny identyfikator.
    \item \texttt{username} (VARCHAR, UNIQUE) -- nazwa wyświetlana.
    \item \texttt{email} (VARCHAR, UNIQUE) -- adres email.
    \item \texttt{password} (VARCHAR) -- skrót hasła (BCrypt).
    \item \texttt{role} (VARCHAR) -- rola systemowa (USER, ADMIN).
    \item \texttt{created\_at} (TIMESTAMP) -- data rejestracji.
\end{itemize}

\subsection{Tabela game\_rooms}

Przechowuje konfigurację i stan pokojów gry:
\begin{itemize}
    \item \texttt{id} (UUID, PK) -- unikalny identyfikator pokoju.
    \item \texttt{name} (VARCHAR) -- nazwa pokoju.
    \item \texttt{room\_code} (VARCHAR, UNIQUE) -- kod dostępu.
    \item \texttt{host\_id} (UUID, FK) -- referencja do twórcy pokoju.
    \item \texttt{status} (VARCHAR) -- aktualny stan (WAITING, IN\_GAME, FINISHED).
    \item \texttt{max\_players} (INTEGER) -- maksymalna liczba graczy.
    \item \texttt{current\_phase} (VARCHAR) -- aktualna faza gry.
\end{itemize}

\subsection{Tabela players\_in\_rooms}

Tabela asocjacyjna łącząca użytkowników z pokojami:
\begin{itemize}
    \item \texttt{id} (UUID, PK) -- unikalny identyfikator rekordu.
    \item \texttt{user\_id} (UUID, FK) -- referencja do użytkownika.
    \item \texttt{room\_id} (UUID, FK) -- referencja do pokoju.
    \item \texttt{game\_role} (VARCHAR) -- przydzielona rola w grze (MAFIA, CITIZEN).
    \item \texttt{is\_alive} (BOOLEAN) -- status gracza w rozgrywce.
    \item \texttt{joined\_at} (TIMESTAMP) -- moment dołączenia.
\end{itemize}

\subsection{Tabela refresh\_tokens}

Obsługuje mechanizm odświeżania tokenów JWT:
\begin{itemize}
    \item \texttt{id} (BIGINT, PK) -- identyfikator.
    \item \texttt{token} (VARCHAR, UNIQUE) -- wartość tokenu.
    \item \texttt{user\_id} (UUID, FK) -- właściciel tokenu.
    \item \texttt{expiry\_date} (TIMESTAMP) -- data wygaśnięcia.
    \item \texttt{revoked} (BOOLEAN) -- flaga unieważnienia.
\end{itemize}
\chapter{Implementacja systemu uwierzytelniania i komunikacji}
\label{chap:implementacja_auth}

Niniejszy rozdział opisuje implementację dwóch kluczowych podsystemów: mechanizmu uwierzytelniania opartego na tokenach JWT oraz warstwy komunikacji w czasie rzeczywistym wykorzystującej protokół WebSocket.

\section{Uwierzytelnianie z wykorzystaniem JWT}

W architekturze aplikacji webowej z oddzielonym frontendem i backendem konieczne jest zastosowanie mechanizmu uwierzytelniania bezstanowego. Wybrano standard JSON Web Token (JWT), który pozwala na weryfikację tożsamości użytkownika bez konieczności przechowywania sesji po stronie serwera.

\subsection{Struktura tokenu JWT}

Token JWT składa się z trzech części oddzielonych kropkami:
\begin{enumerate}
    \item \textbf{Header} -- nagłówek określający algorytm podpisu (HS256).
    \item \textbf{Payload} -- ładunek zawierający dane użytkownika (claims): identyfikator, nazwa, rola, czas wygaśnięcia.
    \item \textbf{Signature} -- podpis cyfrowy generowany z użyciem tajnego klucza serwera.
\end{enumerate}

\subsection{Przepływ uwierzytelniania}

Proces uwierzytelniania przebiega w następujących krokach:
\begin{enumerate}
    \item Klient wysyła żądanie POST na endpoint \texttt{/api/auth/login} z danymi uwierzytelniającymi (email, hasło).
    \item Serwer weryfikuje poświadczenia względem bazy danych (porównanie skrótu BCrypt).
    \item Po pomyślnej weryfikacji serwer generuje parę tokenów: Access Token (krótkotrwały, 15 minut) oraz Refresh Token (długotrwały, 7 dni).
    \item Access Token jest zwracany w ciele odpowiedzi, Refresh Token zapisywany w bazie danych i zwracany jako HttpOnly cookie.
    \item Klient dołącza Access Token do nagłówka \texttt{Authorization: Bearer <token>} przy kolejnych żądaniach.
\end{enumerate}

\subsection{Filtr JWT w Spring Security}

Weryfikacja tokenów została zaimplementowana jako filtr w łańcuchu Spring Security. Klasa \texttt{JwtAuthenticationFilter} rozszerza \texttt{OncePerRequestFilter} i realizuje następującą logikę:

\begin{enumerate}
    \item Ekstrakcja tokenu z nagłówka \texttt{Authorization}.
    \item Walidacja podpisu i daty wygaśnięcia przy użyciu biblioteki JJWT.
    \item Pobranie danych użytkownika (claims) z tokenu.
    \item Utworzenie obiektu \texttt{Authentication} i wstrzyknięcie do \texttt{SecurityContextHolder}.
\end{enumerate}

Filtr jest rejestrowany w konfiguracji \texttt{SecurityFilterChain} przed standardowym filtrem \texttt{UsernamePasswordAuthenticationFilter}.

\subsection{Mechanizm odświeżania tokenów}

Krótki czas życia Access Tokenu (15 minut) wymaga mechanizmu odświeżania bez ponownego logowania. Implementacja opiera się na tabeli \texttt{refresh\_tokens} w bazie danych:

\begin{enumerate}
    \item Klient wykrywa wygaśnięcie Access Tokenu (kod odpowiedzi 401).
    \item Klient wysyła żądanie na endpoint \texttt{/api/auth/refresh} z Refresh Tokenem.
    \item Serwer weryfikuje istnienie i ważność Refresh Tokenu w bazie.
    \item Po pomyślnej weryfikacji generowany jest nowy Access Token.
    \item Opcjonalnie: rotacja Refresh Tokenu (unieważnienie starego, wydanie nowego).
\end{enumerate}

Mechanizm rotacji zapewnia dodatkowe bezpieczeństwo -- wykryty wyciek tokenu skutkuje jego unieważnieniem przy pierwszej próbie użycia.

\section{Konfiguracja Spring Security}

Klasa \texttt{SecurityConfig} definiuje polityki bezpieczeństwa aplikacji poprzez konfigurację \texttt{SecurityFilterChain}.

\subsection{Endpointy publiczne i chronione}

Konfiguracja dzieli endpointy na kategorie:
\begin{itemize}
    \item \textbf{Publiczne} -- dostępne bez uwierzytelnienia: \texttt{/api/auth/login}, \texttt{/api/auth/register}, \texttt{/api/auth/refresh}.
    \item \textbf{Chronione} -- wymagające ważnego tokenu JWT: wszystkie pozostałe endpointy \texttt{/api/**}.
    \item \textbf{Administracyjne} -- wymagające roli ADMIN: \texttt{/api/admin/**}.
\end{itemize}

\subsection{Konfiguracja CORS}

Ze względu na rozdzielenie frontendu (port 3000) i backendu (port 8080) konieczna jest konfiguracja Cross-Origin Resource Sharing:
\begin{itemize}
    \item Dozwolone originy: \texttt{http://localhost:3000}.
    \item Dozwolone metody: GET, POST, PUT, DELETE, OPTIONS.
    \item Dozwolone nagłówki: Authorization, Content-Type.
    \item Credentials: włączone (dla cookies z Refresh Tokenem).
\end{itemize}

\section{Komunikacja w czasie rzeczywistym -- WebSocket}

Tradycyjny model HTTP request-response nie spełnia wymagań aplikacji czasu rzeczywistego. Dla synchronizacji stanu gry między wieloma klientami zastosowano protokół WebSocket z warstwą abstrakcji STOMP.

\subsection{Konfiguracja WebSocket w Spring Boot}

Klasa \texttt{WebSocketConfig} implementuje interfejs \texttt{WebSocketMessageBrokerConfigurer} i definiuje:

\begin{itemize}
    \item \textbf{Endpoint połączenia} -- \texttt{/ws} z obsługą SockJS dla przeglądarek bez natywnego wsparcia WebSocket.
    \item \textbf{Prefiks aplikacji} -- \texttt{/app} dla wiadomości wysyłanych przez klienta do serwera.
    \item \textbf{Prefiks brokera} -- \texttt{/topic} dla kanałów publicznych (broadcast) i \texttt{/queue} dla wiadomości prywatnych.
\end{itemize}

\subsection{Model publikacji i subskrypcji}

Komunikacja opiera się na wzorcu publish-subscribe:
\begin{itemize}
    \item \textbf{Kanały publiczne} (\texttt{/topic/room/\{roomId\}}) -- wszyscy gracze w pokoju subskrybują ten sam kanał i otrzymują aktualizacje stanu gry.
    \item \textbf{Kanały prywatne} (\texttt{/user/queue/private}) -- wiadomości kierowane do konkretnego użytkownika (np. informacja o przydzielonej roli).
\end{itemize}

\subsection{Uwierzytelnianie połączeń WebSocket}

Połączenia WebSocket wymagają autoryzacji w fazie handshake. Implementacja wykorzystuje \texttt{ChannelInterceptor}:
\begin{enumerate}
    \item Klient dołącza token JWT jako parametr zapytania lub nagłówek STOMP.
    \item Interceptor przechwytuje ramkę CONNECT.
    \item Token jest walidowany analogicznie jak w filtrze HTTP.
    \item Po pomyślnej walidacji połączenie jest akceptowane, w przeciwnym razie odrzucane.
\end{enumerate}

\subsection{Obsługa rozłączeń}

System musi reagować na nieoczekiwane rozłączenia graczy (utrata połączenia, zamknięcie przeglądarki). Implementacja \texttt{SessionDisconnectEvent} pozwala na:
\begin{itemize}
    \item Wykrycie rozłączenia użytkownika.
    \item Aktualizację stanu gracza w pokoju (oznaczenie jako nieaktywny).
    \item Powiadomienie pozostałych graczy o rozłączeniu.
    \item Opcjonalnie: automatyczne wstrzymanie gry lub eliminacja gracza po przekroczeniu limitu czasu.
\end{itemize}

\section{Integracja frontendu z mechanizmami bezpieczeństwa}

Warstwa klienta (React) musi prawidłowo obsługiwać cykl życia tokenów i połączenia WebSocket.

\subsection{Przechowywanie tokenów}

Access Token jest przechowywany w pamięci aplikacji (zmienna stanu React) -- nie w localStorage ze względów bezpieczeństwa (podatność na XSS). Refresh Token jest zarządzany jako HttpOnly cookie, niedostępny dla kodu JavaScript.

\subsection{Interceptor Axios}

Biblioteka Axios została skonfigurowana z interceptorami obsługującymi:
\begin{itemize}
    \item Automatyczne dołączanie tokenu do nagłówka Authorization.
    \item Przechwytywanie odpowiedzi 401 i inicjowanie procesu odświeżania tokenu.
    \item Ponawianie oryginalnego żądania po uzyskaniu nowego tokenu.
\end{itemize}

\subsection{Klient STOMP}

Połączenie WebSocket jest nawiązywane przy użyciu biblioteki \texttt{@stomp/stompjs} z konfiguracją:
\begin{itemize}
    \item Automatyczne ponowne łączenie przy utracie połączenia.
    \item Przekazywanie tokenu JWT w nagłówku CONNECT.
    \item Obsługa błędów połączenia z informowaniem użytkownika.
\end{itemize}
\chapter{Implementacja logiki rozgrywki}
\label{chap:implementacja_rozgrywki}

Niniejszy rozdział opisuje implementację rdzenia systemu -- silnika gry odpowiedzialnego za zarządzanie fazami rozgrywki, przetwarzanie głosowań i weryfikację warunków zwycięstwa.

\section{Architektura silnika gry}

Silnik gry został zaprojektowany z wykorzystaniem wzorca maszyny stanowej (State Machine), który naturalnie odwzorowuje cykliczny charakter rozgrywki \textit{Mafia}.

\subsection{Klasa GamePhaseStateMachine}

Centralnym elementem silnika jest klasa \texttt{GamePhaseStateMachine}, która:
\begin{itemize}
    \item Przechowuje definicje wszystkich możliwych stanów gry.
    \item Definiuje dozwolone przejścia między stanami.
    \item Zarządza aktualnym stanem pokoju gry.
    \item Wyzwala akcje przy wejściu i wyjściu ze stanu.
\end{itemize}

\subsection{Stany rozgrywki}

System definiuje następujące stany reprezentowane przez enum \texttt{GamePhase}:

\begin{itemize}
    \item \texttt{LOBBY} -- stan początkowy, gracze dołączają do pokoju, host konfiguruje ustawienia.
    \item \texttt{NIGHT\_VOTING} -- faza nocna, członkowie Mafii wybierają ofiarę.
    \item \texttt{NIGHT\_RESULT} -- prezentacja wyniku nocy, informacja o eliminacji lub jej braku.
    \item \texttt{DAY\_VOTING} -- faza dzienna, wszyscy żywi gracze głosują na podejrzanego.
    \item \texttt{DAY\_RESULT} -- prezentacja wyniku głosowania dziennego.
    \item \texttt{GAME\_OVER} -- stan końcowy, wyświetlenie zwycięskiej frakcji i statystyk.
\end{itemize}

\subsection{Przejścia między stanami}

Przejścia są definiowane w metodzie \texttt{defineTransitions()} i wyzwalane przez:
\begin{itemize}
    \item Akcje użytkownika (host rozpoczyna grę, gracze oddają głosy).
    \item Zdarzenia systemowe (upłynął czas fazy, zebrano wszystkie głosy).
    \item Warunki logiczne (spełniony warunek zwycięstwa).
\end{itemize}

Przykładowy cykl przejść:
\begin{center}
LOBBY $\rightarrow$ NIGHT\_VOTING $\rightarrow$ NIGHT\_RESULT $\rightarrow$ DAY\_VOTING $\rightarrow$ DAY\_RESULT $\rightarrow$ NIGHT\_VOTING $\rightarrow$ ...
\end{center}

Cykl jest przerywany przejściem do \texttt{GAME\_OVER} gdy zostanie spełniony warunek zwycięstwa.

\section{System głosowania}

Głosowanie stanowi główny mechanizm interakcji graczy z systemem. Implementacja musi obsługiwać dwa różne konteksty: głosowanie nocne (Mafia) i głosowanie dzienne (wszyscy gracze).

\subsection{Klasa VotingSessionService}

Serwis \texttt{VotingSessionService} zarządza cyklem życia sesji głosowania:
\begin{enumerate}
    \item \textbf{Inicjalizacja} -- utworzenie nowej sesji z listą uprawnionych głosujących.
    \item \textbf{Zbieranie głosów} -- rejestracja głosów z walidacją uprawnień.
    \item \textbf{Obliczanie wyniku} -- zliczenie głosów i wyłonienie wyniku.
    \item \textbf{Rozstrzyganie remisów} -- zastosowanie reguły rozstrzygającej (brak eliminacji lub losowanie).
    \item \textbf{Zamknięcie} -- zakończenie sesji i publikacja wyniku.
\end{enumerate}

\subsection{Walidacja głosów}

Każdy oddany głos przechodzi walidację:
\begin{itemize}
    \item Czy głosujący jest uczestnikiem danego pokoju?
    \item Czy głosujący jest żywy (nie został wyeliminowany)?
    \item Czy głosujący ma prawo głosu w bieżącej fazie (np. tylko Mafia w nocy)?
    \item Czy cel głosu jest prawidłowy (żywy gracz, nie samogłos jeśli zabroniony)?
    \item Czy głosujący nie oddał już głosu w tej sesji?
\end{itemize}

Nieprawidłowe głosy są odrzucane z odpowiednim komunikatem błędu.

\subsection{Algorytm zliczania głosów}

Po zebraniu wszystkich głosów lub upływie czasu system wykonuje:
\begin{enumerate}
    \item Grupowanie głosów według celu.
    \item Sortowanie celów według liczby otrzymanych głosów.
    \item Sprawdzenie czy istnieje jednoznaczny zwycięzca.
    \item W przypadku remisu -- zastosowanie reguły konfigurowalnej (brak eliminacji / losowanie).
    \item Publikacja wyniku do wszystkich klientów.
\end{enumerate}

\section{Zarządzanie eliminacjami}

Eliminacja gracza jest konsekwencją wyniku głosowania i wpływa na stan całej rozgrywki.

\subsection{Proces eliminacji}

Po ustaleniu wyniku głosowania system:
\begin{enumerate}
    \item Aktualizuje status gracza w bazie danych (\texttt{is\_alive = false}).
    \item Publikuje zdarzenie eliminacji do kanału pokoju.
    \item Ujawnia rolę wyeliminowanego gracza wszystkim uczestnikom.
    \item Wyzwala weryfikację warunków zwycięstwa.
\end{enumerate}

\subsection{Weryfikacja warunków zwycięstwa}

Klasa \texttt{GameOrchestrator} implementuje metodę \texttt{checkWinConditions()}:

\begin{itemize}
    \item Pobiera liczbę żywych graczy z rolą MAFIA.
    \item Pobiera liczbę żywych graczy z rolą CITIZEN.
    \item Jeśli liczba Mafii $\geq$ liczba Obywateli -- zwycięstwo Mafii.
    \item Jeśli liczba Mafii = 0 -- zwycięstwo Obywateli.
    \item W przeciwnym razie -- gra kontynuowana.
\end{itemize}

\section{Przydzielanie ról}

Na początku rozgrywki system musi losowo przydzielić role zachowując odpowiednie proporcje.

\subsection{Algorytm przydziału}

\begin{enumerate}
    \item Pobranie listy graczy w pokoju.
    \item Obliczenie liczby członków Mafii na podstawie konfiguracji (domyślnie: 1 Mafia na każdych 4 graczy).
    \item Losowe wybranie graczy do frakcji Mafia.
    \item Przypisanie pozostałym graczom roli Citizen.
    \item Zapisanie ról w bazie danych.
    \item Wysłanie informacji o roli do każdego gracza przez kanał prywatny.
\end{enumerate}

\subsection{Bezpieczeństwo informacji o rolach}

Role są przechowywane po stronie serwera i nigdy nie są ujawniane w komunikacji publicznej przed zakończeniem gry. Każdy gracz otrzymuje informację wyłącznie o swojej roli przez dedykowany kanał WebSocket (\texttt{/user/queue/role}).

\section{Orkiestracja rozgrywki}

Klasa \texttt{GameOrchestrator} koordynuje wszystkie komponenty silnika gry.

\subsection{Rozpoczęcie gry}

Metoda \texttt{startGame()} wykonuje sekwencję:
\begin{enumerate}
    \item Walidacja minimalnej liczby graczy.
    \item Zmiana statusu pokoju na IN\_GAME.
    \item Uruchomienie algorytmu przydziału ról.
    \item Przejście maszyny stanowej do NIGHT\_VOTING.
    \item Uruchomienie timera fazy nocnej.
    \item Publikacja stanu gry do wszystkich klientów.
\end{enumerate}

\subsection{Obsługa zakończenia fazy}

Timer fazy wyzwala zdarzenie zakończenia. System:
\begin{enumerate}
    \item Zamyka aktywną sesję głosowania (jeśli istnieje).
    \item Oblicza wynik na podstawie zebranych głosów.
    \item Wykonuje eliminację (jeśli dotyczy).
    \item Sprawdza warunki zwycięstwa.
    \item Przechodzi do następnej fazy lub kończy grę.
\end{enumerate}

\section{Implementacja warstwy klienta}

Frontend React musi efektywnie reagować na zmiany stanu gry otrzymywane przez WebSocket.

\subsection{Struktura komponentów gry}

Widok gry składa się z hierarchii komponentów:
\begin{itemize}
    \item \texttt{GameView} -- kontener główny, zarządza połączeniem WebSocket.
    \item \texttt{GamePanel} -- wyświetla aktualną fazę i timer.
    \item \texttt{PlayerList} -- lista graczy z oznaczeniem statusu (żywy/martwy).
    \item \texttt{VotingPanel} -- interfejs głosowania (widoczny w odpowiednich fazach).
    \item \texttt{GameResultPanel} -- ekran końcowy z wynikiem i statystykami.
\end{itemize}

\subsection{Zarządzanie stanem w React}

Stan gry jest przechowywany w komponencie \texttt{GameView} i propagowany do komponentów potomnych przez props. Aktualizacje stanu są wyzwalane przez:
\begin{itemize}
    \item Wiadomości WebSocket z kanału \texttt{/topic/room/\{id\}/game-state}.
    \item Prywatne wiadomości z kanału \texttt{/user/queue/private}.
    \item Lokalne akcje użytkownika (oddanie głosu).
\end{itemize}

\subsection{Obsługa faz po stronie klienta}

Komponent renderuje różne interfejsy w zależności od aktualnej fazy:
\begin{itemize}
    \item \texttt{NIGHT\_VOTING} -- Mafia widzi panel wyboru ofiary, Obywatele widzą ekran oczekiwania.
    \item \texttt{DAY\_VOTING} -- wszyscy żywi gracze widzą panel głosowania.
    \item \texttt{*\_RESULT} -- wszyscy widzą wynik i oczekują na przejście.
    \item \texttt{GAME\_OVER} -- panel końcowy z ujawnionymi rolami.
\end{itemize}
\chapter{Wdrożenie i testowanie}
\label{chap:wdrozenie}

Niniejszy rozdział opisuje infrastrukturę wdrożeniową projektu, strategię konteneryzacji oraz podejście do testowania systemu.

\section{Konteneryzacja z Docker}

Cała infrastruktura aplikacji została zdefiniowana jako kod (Infrastructure as Code) w pliku \texttt{docker-compose.yml}. Podejście to zapewnia powtarzalność środowiska i eliminuje problemy typu \textit{works on my machine}.

\subsection{Architektura kontenerów}

System składa się z pięciu kontenerów Docker:

\begin{enumerate}
    \item \textbf{mafia-backend} -- aplikacja Spring Boot (port 8080)
    \begin{itemize}
        \item Obraz budowany z \texttt{backend/Dockerfile}.
        \item Zależny od kontenerów \texttt{db} i \texttt{rabbitmq}.
        \item Zmienne środowiskowe konfigurują połączenia z bazą i brokerem.
    \end{itemize}
    
    \item \textbf{mafia-frontend} -- aplikacja React serwowana przez Nginx (port 3000)
    \begin{itemize}
        \item Obraz budowany z \texttt{frontend/Dockerfile}.
        \item Wieloetapowy build: Node.js do kompilacji, Nginx do serwowania.
        \item Konfiguracja proxy dla żądań API.
    \end{itemize}
    
    \item \textbf{db} -- baza danych PostgreSQL 15 (port 5432)
    \begin{itemize}
        \item Obraz: \texttt{postgres:15-alpine}.
        \item Wolumen \texttt{postgres\_data} dla trwałości danych.
        \item Healthcheck weryfikujący gotowość bazy.
    \end{itemize}
    
    \item \textbf{rabbitmq} -- broker wiadomości RabbitMQ 3.13 (porty 5672, 15672)
    \begin{itemize}
        \item Obraz: \texttt{rabbitmq:3.13-management-alpine}.
        \item Panel zarządzania dostępny na porcie 15672.
        \item Healthcheck sprawdzający status brokera.
    \end{itemize}
    
    \item \textbf{mafia-backend-tests} -- kontener testowy
    \begin{itemize}
        \item Uruchamiany z profilem Spring \texttt{test}.
        \item Wykonuje testy przy starcie (\texttt{mvn test}).
        \item Wyniki zapisywane do wolumenu \texttt{test-reports}.
    \end{itemize}
\end{enumerate}

\subsection{Mechanizm healthcheck}

Kontenery infrastrukturalne (PostgreSQL, RabbitMQ) posiadają zdefiniowane healthchecki:

\begin{itemize}
    \item PostgreSQL: \texttt{pg\_isready -U postgres}
    \item RabbitMQ: \texttt{rabbitmq-diagnostics check\_running}
\end{itemize}

Kontener \texttt{mafia-backend} deklaruje zależność \texttt{depends\_on} z warunkiem \texttt{service\_healthy}, co gwarantuje uruchomienie dopiero po pełnej inicjalizacji infrastruktury.

\subsection{Dockerfile backendu}

Dockerfile backendu wykorzystuje wieloetapowy build:
\begin{enumerate}
    \item \textbf{Etap build} -- obraz Maven, kompilacja projektu, generowanie JAR.
    \item \textbf{Etap runtime} -- lekki obraz JRE, kopiowanie artefaktu, ekspozycja portu 8080.
\end{enumerate}

Podejście to minimalizuje rozmiar finalnego obrazu (brak narzędzi buildowych w produkcji).

\subsection{Dockerfile frontendu}

Frontend również wykorzystuje wieloetapowy build:
\begin{enumerate}
    \item \textbf{Etap build} -- obraz Node.js, instalacja zależności, \texttt{npm run build}.
    \item \textbf{Etap runtime} -- obraz Nginx, kopiowanie zbudowanych plików do \texttt{/usr/share/nginx/html}.
\end{enumerate}

Konfiguracja Nginx (\texttt{nginx.conf}) definiuje proxy dla żądań \texttt{/api/*} kierowanych do backendu.

\section{Uruchamianie środowiska}

\subsection{Wymagania}

Do uruchomienia projektu wymagane są:
\begin{itemize}
    \item Docker Engine w wersji 20.10 lub nowszej.
    \item Docker Compose w wersji 2.0 lub nowszej.
    \item Minimum 4 GB RAM dostępnego dla kontenerów.
\end{itemize}

\subsection{Procedura uruchomienia}

\begin{enumerate}
    \item Sklonowanie repozytorium: \texttt{git clone <url>}
    \item Przejście do katalogu projektu: \texttt{cd eMafia}
    \item Skopiowanie pliku konfiguracyjnego: \texttt{cp exampleEnvFile .env}
    \item Uruchomienie kontenerów: \texttt{docker-compose up -d}
    \item Weryfikacja statusu: \texttt{docker-compose ps}
\end{enumerate}

Aplikacja jest dostępna pod adresem \texttt{http://localhost:3000}.

\section{Testowanie}

Strategia testowania obejmuje testy jednostkowe i integracyjne wykonywane w izolowanym środowisku kontenerowym.

\subsection{Testy jednostkowe}

Testy jednostkowe weryfikują izolowane fragmenty logiki biznesowej:
\begin{itemize}
    \item Algorytmy przydziału ról.
    \item Logika zliczania głosów i rozstrzygania remisów.
    \item Walidacja danych wejściowych.
    \item Generowanie i walidacja tokenów JWT.
\end{itemize}

Framework testowy: JUnit 5 z biblioteką Mockito do izolacji zależności.

\subsection{Testy integracyjne}

Testy integracyjne weryfikują współpracę komponentów:
\begin{itemize}
    \item Komunikacja z bazą danych (Spring Data JPA).
    \item Endpointy REST API (MockMvc).
    \item Komunikacja WebSocket (testowe klienty STOMP).
\end{itemize}

Profil testowy (\texttt{application-test.properties}) konfiguruje bazę H2 w pamięci, przyspieszając wykonanie testów.

\subsection{Uruchamianie testów}

Testy można uruchomić na dwa sposoby:
\begin{enumerate}
    \item Lokalnie: \texttt{cd backend \&\& mvn test}
    \item W kontenerze: \texttt{docker-compose run mafia-backend-tests}
\end{enumerate}

Raporty testów są generowane w formacie Surefire i zapisywane do katalogu \texttt{target/surefire-reports}.

\section{Monitorowanie i logi}

\subsection{Logowanie aplikacji}

Backend wykorzystuje framework Logback z konfiguracją w pliku \texttt{logback-spring.xml}:
\begin{itemize}
    \item Poziomy logowania konfigurowalne per pakiet.
    \item Logi zapisywane do plików z rotacją dzienną.
    \item Format logu zawierający timestamp, poziom, klasę i wiadomość.
\end{itemize}

\subsection{Dostęp do logów kontenerów}

Logi kontenerów są dostępne przez Docker:
\begin{itemize}
    \item Podgląd na żywo: \texttt{docker-compose logs -f mafia-backend}
    \item Logi konkretnego kontenera: \texttt{docker logs <container\_id>}
\end{itemize}

\subsection{Panel RabbitMQ}

Panel zarządzania RabbitMQ (port 15672) umożliwia:
\begin{itemize}
    \item Monitorowanie kolejek i wymian.
    \item Podgląd statystyk wiadomości.
    \item Ręczne publikowanie wiadomości testowych.
\end{itemize}

Domyślne poświadczenia: guest/guest.

\section{Możliwości rozwoju}

\subsection{Rozszerzenie systemu ról}

Aktualna implementacja obsługuje dwie role podstawowe (Mafia, Obywatel). System został zaprojektowany z myślą o rozszerzeniu o role specjalne:
\begin{itemize}
    \item Detektyw -- możliwość sprawdzenia roli wybranego gracza w nocy.
    \item Lekarz -- ochrona wybranego gracza przed eliminacją nocną.
    \item Strażnik -- blokowanie akcji wybranego gracza.
\end{itemize}

Dodanie nowej roli wymaga:
\begin{enumerate}
    \item Rozszerzenia enumu \texttt{GameRole}.
    \item Implementacji logiki akcji nocnej w \texttt{NightPhaseHandler}.
    \item Aktualizacji algorytmu przydziału ról.
    \item Dodania odpowiedniego interfejsu w komponencie React.
\end{enumerate}

\subsection{Wdrożenie produkcyjne}

Dla wdrożenia produkcyjnego zalecane są następujące modyfikacje:
\begin{itemize}
    \item Konfiguracja HTTPS z certyfikatami SSL.
    \item Zewnętrzna baza danych (managed PostgreSQL).
    \item Skalowanie horyzontalne backendu za load balancerem.
    \item Konfiguracja RabbitMQ w trybie klastra.
    \item Monitoring z Prometheus i Grafana.
\end{itemize}
\chapter*{Podsumowanie}
\addcontentsline{toc}{chapter}{Podsumowanie}

Niniejsza praca inżynierska przedstawiła proces projektowania i implementacji rozproszonego systemu do prowadzenia gry towarzyskiej \textit{Mafia} w środowisku przeglądarkowym. Zrealizowany projekt stanowi odpowiedź na zidentyfikowaną lukę technologiczną -- brak platformy łączącej automatyzację logiki gry z zachowaniem bezpośredniej interakcji między uczestnikami.

\section*{Realizacja celów}

Główny cel pracy -- stworzenie działającego prototypu platformy \textit{eMafia} -- został osiągnięty. System realizuje następujące założenia:

\begin{itemize}
    \item \textbf{Automatyzacja funkcji Mistrza Gry} -- silnik gry samodzielnie zarządza fazami rozgrywki, głosowaniami i eliminacjami, eliminując potrzebę obecności prowadzącego.
    
    \item \textbf{Komunikacja w czasie rzeczywistym} -- wykorzystanie protokołu WebSocket z warstwą STOMP zapewnia natychmiastową synchronizację stanu gry między wszystkimi klientami.
    
    \item \textbf{Bezpieczeństwo} -- mechanizm uwierzytelniania JWT chroni dostęp do zasobów, a architektura systemu zapewnia poufność informacji o rolach graczy.
    
    \item \textbf{Konteneryzacja} -- pełna definicja infrastruktury w Docker Compose umożliwia powtarzalne wdrożenie środowiska.
\end{itemize}

\section*{Zastosowane technologie}

Projekt wykorzystuje nowoczesny stos technologiczny:
\begin{itemize}
    \item Backend: Java 21, Spring Boot 3.3, Spring Security, Spring WebSocket.
    \item Frontend: React 18, CSS Modules, Axios, SockJS/STOMP.
    \item Infrastruktura: Docker, Docker Compose, PostgreSQL 15, RabbitMQ 3.13.
\end{itemize}

Wybór technologii był podyktowany wymaganiami projektu -- koniecznością obsługi komunikacji dwukierunkowej, bezpieczeństwa danych i łatwości wdrożenia.

\section*{Wnioski}

Realizacja projektu potwierdziła, że:
\begin{enumerate}
    \item Wzorzec maszyny stanowej jest naturalnym i efektywnym modelem dla logiki gry turowej.
    \item Protokół WebSocket z warstwą STOMP doskonale sprawdza się w aplikacjach wymagających synchronizacji stanu w czasie rzeczywistym.
    \item Konteneryzacja znacząco upraszcza proces wdrożenia i eliminuje problemy środowiskowe.
\end{enumerate}

\section*{Kierunki dalszego rozwoju}

Zrealizowany prototyp stanowi fundament dla dalszych rozszerzeń:
\begin{itemize}
    \item Dodanie ról specjalnych (Detektyw, Lekarz, Strażnik) zwiększy głębię strategiczną rozgrywki.
    \item Implementacja systemu czatu w grze umożliwi komunikację tekstową między graczami.
    \item Wdrożenie produkcyjne w środowisku chmurowym (np. AWS, Azure) pozwoli na publiczne udostępnienie platformy.
    \item Rozwój aplikacji mobilnej (React Native) poprawi dostępność dla użytkowników mobilnych.
\end{itemize}

\chapter*{Załączniki}
\addcontentsline{toc}{chapter}{Załączniki}

\section*{Struktura projektu}

Repozytorium projektu zawiera następujące katalogi:
\begin{itemize}
    \item \texttt{backend/} -- kod źródłowy aplikacji Spring Boot.
    \item \texttt{frontend/} -- kod źródłowy aplikacji React.
    \item \texttt{docker-compose.yml} -- definicja infrastruktury kontenerowej.
    \item \texttt{example\_database\_data.sql} -- przykładowe dane testowe.
    \item \texttt{exampleEnvFile} -- szablon zmiennych środowiskowych.
\end{itemize}

\section*{Instrukcja uruchomienia}

Szczegółowa instrukcja uruchomienia projektu znajduje się w pliku \texttt{README.md} w głównym katalogu repozytorium.

\begin{thebibliography}{99}

\bibitem{martin2008}
Martin, R. C. (2008). \textit{Clean Code: A Handbook of Agile Software Craftsmanship}. Prentice Hall.

\bibitem{spring2024}
Spring Framework Documentation. (2024). \textit{Spring Boot Reference Documentation}. https://docs.spring.io/spring-boot/

\bibitem{websocket2011}
Fette, I., Melnikov, A. (2011). \textit{The WebSocket Protocol}. RFC 6455. IETF.

\bibitem{stomp2012}
STOMP Protocol Specification. (2012). \textit{Simple Text Oriented Messaging Protocol}. https://stomp.github.io/

\bibitem{jwt2015}
Jones, M., Bradley, J., Sakimura, N. (2015). \textit{JSON Web Token (JWT)}. RFC 7519. IETF.

\bibitem{docker2024}
Docker Documentation. (2024). \textit{Docker Compose Reference}. https://docs.docker.com/compose/

\bibitem{react2024}
React Documentation. (2024). \textit{React: A JavaScript library for building user interfaces}. https://react.dev/

\bibitem{gamma1994}
Gamma, E., Helm, R., Johnson, R., Vlissides, J. (1994). \textit{Design Patterns: Elements of Reusable Object-Oriented Software}. Addison-Wesley.

\bibitem{postgresql2024}
PostgreSQL Global Development Group. (2024). \textit{PostgreSQL 15 Documentation}. https://www.postgresql.org/docs/15/

\bibitem{rabbitmq2024}
RabbitMQ Documentation. (2024). \textit{RabbitMQ Tutorials}. https://www.rabbitmq.com/getstarted.html

\end{thebibliography}

\end{document}